% Chapter 5 - The Interactive Unity Application

\section{The Interactive Unity Application}\label{sec:application}

The pipeline of the previous chapter ends at a socket. For every frame in
which at least one hand is visible, the streamer sends a \acrshort{json}
datagram to the loopback address on port \texttt{5052}, carrying a timestamp
and, for each hand, the sixteen axis-angle joint rotations, a gesture label,
and a placement anchor~(\autoref{subsec:streaming},~\autoref{fig:udp-schema}).
This chapter begins where that datagram arrives. It describes the Unity
application that consumes the stream: the receiver that decodes it into a live
skeleton, the interaction layer that every scene reads from, and the four
scenes that put the recovered pose to different uses.

The presentation moves outward from the stream. The receiver is treated first,
as the exact inverse of the serialisation step that closed the pipeline. The
shared interaction layer follows, since the gesture vocabulary, the scene
navigation, and the body-anchored menu it provides recur across every scene and
are best described once. The four scenes are then taken in turn. Each is
examined not only as a mechanism but against the state of the
art~(\autoref{subsec:sota_interaction}): the question throughout is what input
hardware a comparable system would demand, and what the present scene achieves
without it.

A single architectural principle organises the whole application. The receiver
is the only component that touches the network, and it performs no interaction
logic of its own. Every scene component holds a reference to it and reads the
current skeleton, gesture, and anchor each frame. The reconstructed hand is thus
a shared resource, polled by many independent consumers, and the scenes remain
entirely decoupled from the transport beneath them.
